\section{Appendix A}
\subsection{Formulas}
This section summarizes the objectives implemented in the planner using notation close to the original ALM latent-collocation formulation. Let
\(\mathcal{Z}=\{z_0,\ldots,z_T\}\) denote the optimized latent trajectory and
\(\mathcal{A}=\{a_0,\ldots,a_{T-1}\}\) denote the optimized macro-action sequence. Each \(a_t\in\mathbb{R}^{d_a}\) is a macro-action. In Push-T, \(d_a=10\), corresponding to five consecutive 2D primitive actions packed into one action vector:
\[
a_t = [u_{5t},u_{5t+1},u_{5t+2},u_{5t+3},u_{5t+4}],
\qquad
u_i\in\mathbb{R}^2.
\]
We write \(h_t\) for the history window used by the world model. Depending on the model, this may contain only latent history \(z_{t-H+1:t}\), or both latent and action history. The pretrained DINO-WM dynamics are denoted by \(f_\psi\).

\paragraph{Base latent task cost.}
Following the ALM formulation, we define a latent-space task cost
\[
min_{\mathcal{Z},\mathcal{A}}
\lambda_v
\sum_{t=1}^{T}
\mathcal{L}_{\mathrm{vid}}(z_t,z_t^{\mathrm{vid}})
+
\lambda_g
\mathcal{L}_{\mathrm{goal}}(z_T,z_g)
+
\lambda_r
\sum_{t=0}^{T-1}
\|a_t\|_2^2 .
\]
Here, the first term aligns the optimized latent trajectory with the generated video plan, the second term penalizes terminal goal error, and the final term regularizes action magnitude.

\paragraph{Original ALM dynamics objective.}
The original latent-collocation planner treats both \(\mathcal{Z}\) and \(\mathcal{A}\) as optimization variables and enforces consistency with the pretrained world model through the dynamics violation
\[
\mathcal{L}^t_{dyn}(\mathcal{Z},\mathcal{A};f_\psi )
=
z_{t+1}-f_\psi(h_t,a_t).
\]
The augmented Lagrangian objective is
\[
\mathcal{L}_{\rho}(\mathcal{Z},\mathcal{A},\Lambda)
=
\tilde C(\mathcal{Z},\mathcal{A})
+
\sum_{t=0}^{T-1}
\left[
\lambda_t^\top \mathcal{L}^t_{dyn}(\mathcal{Z},\mathcal{A})
+
\frac{\rho}{2}
\|\mathcal{L}^t_{dyn}||_2^2
\right],
\]
where \(\Lambda\) are the dual variables and \(\rho\) is the quadratic penalty weight. After each outer ALM iteration, we update
\[
\lambda_t \leftarrow \lambda_t+\rho \mathcal{L}_{dyn}(\mathcal{Z},\mathcal{A}),
\qquad
\rho \leftarrow \min(\gamma\rho,\rho_{\max}).
\]

\paragraph{Soft DINO-WM dynamics.}
As a simpler baseline, we remove the dual variables and use the DINO-WM residual as a soft penalty:
\[
\mathcal{L}_{\mathrm{soft}}(\mathcal{Z},\mathcal{A})
=
\tilde C(\mathcal{Z},\mathcal{A})
+
\lambda_{\mathrm{dyn}}
\sum_{t=0}^{T-1}
\|z_{t+1}-f_\psi(h_t,a_t)\|_2^2 .
\]
This keeps the same collocation variables \((\mathcal{Z},\mathcal{A})\), but replaces the ALM constraint update by a fixed-weight dynamics penalty.

\paragraph{DSM feasibility model.}
We train a conditional denoising model \(\epsilon_\phi\) in the DINO-WM latent space. For a dataset transition \((h_t,a_t,z_{t+1})\), we sample
\[
\epsilon\sim\mathcal{N}(0,I),
\qquad
\sigma\sim p(\sigma),
\]
and corrupt the next latent as
\[
\tilde z_{t+1}=z_{t+1}+\sigma\epsilon.
\]
The DSM objective is the mean-squared denoising error over the latent dimensions:
\[
\mathcal{L}_{\mathrm{DSM}}(\phi)
=
\mathbb{E}_{t,\sigma,\epsilon}
\left[
\frac{1}{D}
\left\|
\epsilon_\phi(\tilde z_{t+1},h_t,a_t,\sigma)-\epsilon
\right\|_2^2
\right],
\]
where \(D\) is the dimensionality of the DINO-WM latent. At planning time, the trained denoising model defines a feasibility energy at a fixed noise level \(\sigma^\ast\):
\[
\mathcal{P}_{\mathrm{DSM}}(z_{t+1},h_t,a_t)
=
\frac{1}{D}
\left\|
\epsilon_\phi(z_{t+1},h_t,a_t,\sigma^\ast)
\right\|_2^2 .
\]

\paragraph{DSM feasibility planning.}
The feasibility-only variant removes the DINO-WM dynamics constraint and augments the task cost with the learned DSM energy:
\[
\mathcal{L}_{\mathrm{feas}}(\mathcal{Z},\mathcal{A})
=
\tilde C(\mathcal{Z},\mathcal{A})
+
\lambda_{\mathrm{feas}}
\frac{1}{T}
\sum_{t=0}^{T-1}
\mathcal{P}_{\mathrm{DSM}}(z_{t+1},h_t,a_t).
\]
Unlike the ALM objective, this formulation does not enforce
\[
z_{t+1}=f_\psi(h_t,a_t),
\]
so the latent trajectory remains a free optimization variable.

\paragraph{Anti-stillness diagnostic.}
To test whether feasibility-only planning collapses to small actions, we add an anti-stillness term
\[
\mathcal{L}_{\mathrm{anti}}(\mathcal{A})
=
\lambda_{\mathrm{anti}}
\frac{1}{T}
\sum_{t=0}^{T-1}
\max(0,m-\|a_t\|_2)^2 .
\]
The resulting objective is
\[
\mathcal{L}_{\mathrm{feas+anti}}
=
\mathcal{L}_{\mathrm{feas}}
+
\mathcal{L}_{\mathrm{anti}}.
\]
This can be seen as the opposite of the last term in the base task where action magnitude is penalized. 
\paragraph{Learned transition head: current-delta target.}
In addition to the DSM head, we train a transition head \(\Delta_\phi\). In the current-delta variant, the target is the latent displacement
\[
\Delta_t = z_{t+1}-z_t.
\]
The training loss is
\[
\mathcal{L}_{\Delta}(\phi)
=
\mathbb{E}_{t}
\left[
\left\|
\Delta_\phi(h_t,a_t)-(z_{t+1}-z_t)
\right\|_2^2
\right].
\]
At planning time, this gives the transition energy
\[
\mathcal{P}_{\Delta}(z_{t+1},h_t,a_t)
=
\left\|
z_{t+1}
-
\left(z_t+\Delta_\phi(h_t,a_t)\right)
\right\|_2^2 .
\]

\paragraph{Learned transition head: paired DINO-WM residual target.}
We also consider a paired residual target relative to the pretrained world model:
\[
r_t^\psi
=
z_{t+1}-f_\psi(h_t,a_t).
\]
The residual head \(r_\phi\) is trained with
\[
\mathcal{L}_{\mathrm{res}}(\phi)
=
\mathbb{E}_{t}
\left[
\left\|
r_\phi(h_t,a_t)
-
\left(z_{t+1}-f_\psi(h_t,a_t)\right)
\right\|_2^2
\right].
\]
The corresponding planning-time transition energy is
\[
\mathcal{P}_{\mathrm{res}}(z_{t+1},h_t,a_t)
=
\left\|
z_{t+1}
-
\left(f_\psi(h_t,a_t)+r_\phi(h_t,a_t)\right)
\right\|_2^2 .
\]

\paragraph{Combined dynamics and feasibility objective.}
The learned feasibility terms can be added to either ALM or soft dynamics. In the soft-dynamics case, the combined objective is
\[
\mathcal{L}_{\mathrm{soft+feas}}
=
\tilde C(\mathcal{Z},\mathcal{A})
+
\lambda_{\mathrm{dyn}}
\sum_{t=0}^{T-1}
\|z_{t+1}-f_\psi(h_t,a_t)\|_2^2
+
\lambda_{\mathrm{feas}}
\sum_{t=0}^{T-1}
\mathcal{P}_{\mathrm{DSM}}(z_{t+1},h_t,a_t)
+
\lambda_{\mathrm{transition}}
\sum_{t=0}^{T-1}
\mathcal{P}_{\mathrm{trans}}(z_{t+1},h_t,a_t),
\]
where \(\mathcal{P}_{\mathrm{trans}}\) is either
\(\mathcal{P}_{\Delta}\) for the current-delta model or
\(\mathcal{P}_{\mathrm{res}}\) for the paired residual model. Setting
\(\lambda_{\mathrm{transition}}=0\) gives DSM-only feasibility as an auxillary on the soft dynamics.

\paragraph{DINO-WM rollout planning.}
In rollout mode, we no longer optimize \(\mathcal{Z}\) directly. Instead, the planner optimizes only the action sequence \(\mathcal{A}\), and the latent trajectory is generated recursively by DINO-WM:
\[
z_{t+1}=f_\psi(h_t,a_t).
\]
Thus, dynamics consistency is enforced by construction rather than by an ALM residual. The rollout objective is
\[
\mathcal{L}_{\mathrm{rollout}}
=
\tilde C(\mathcal{Z}(\mathcal{A}),\mathcal{A})
+
\lambda_{\mathrm{feas}}
\sum_{t=0}^{T-1}
\mathcal{P}_{\mathrm{DSM}}(z_{t+1},h_t,a_t)
+
\lambda_{\mathrm{transition}}
\sum_{t=0}^{T-1}
\mathcal{P}_{\mathrm{trans}}(z_{t+1},h_t,a_t).
\]
For pure DSM rollout, we set
\[
\lambda_{\mathrm{transition}}=0,
\qquad
\lambda_{\mathrm{dyn}}=0,
\]


\paragraph{Action bounds and reparameterization.}
All planner variants optimize bounded normalized DINO-WM actions. When action reparameterization is enabled, we optimize an unconstrained variable \(u_t\) and map it to the valid action range:
\[
a_t
=
a_{\min}
+
\frac{1}{2}
\left(\tanh(u_t)+1\right)
(a_{\max}-a_{\min}).
\]
When reparameterization is disabled, \(a_t\) is optimized directly and clamped to
\([a_{\min},a_{\max}]\).

\paragraph{Expert-action prior diagnostic.}
For oracle warm-start experiments, we additionally penalize deviation from a reference action sequence \(\mathcal{A}^{\mathrm{prior}}\):
\[
\mathcal{L}_{\mathrm{prior}}
=
\lambda_{\mathrm{prior}}
\frac{1}{T}
\sum_{t=0}^{T-1}
\|a_t-a_t^{\mathrm{prior}}\|_2^2 .
\]
This term is used only as a diagnostic to test whether rollout optimization can remain in the expert-action basin.


\section{have done}
\subsection{Soft DINO-WM Dynamics Ablation}
\label{app:soft_dynamics_ablation}

We evaluate a soft DINO-WM dynamics baseline on Push-T to isolate the effect of the pretrained world-model residual and the action regularization weight. In these runs, the dynamics constraint is not handled with ALM dual updates; instead, we use a fixed soft penalty
\[
\mathcal{L}_{\mathrm{soft}}
=
\tilde C(\mathcal{Z},\mathcal{A})
+
\lambda_{\mathrm{dyn}}
\sum_{t=0}^{T-1}
\left\|
z_{t+1}-f_\psi(h_t,a_t)
\right\|_2^2,
\]
with \(\lambda_{\mathrm{dyn}}=10\). The learned feasibility model is disabled in this ablation.

\begin{table}[h]
\centering
\small
\caption{Soft DINO-WM dynamics ablation on Push-T. The results show that the soft world-model residual can produce successful behavior, but the task is sensitive to the action regularization weight.}
\label{tab:soft_dynamics_ablation_appendix}
\begin{tabular}{lccccp{0.28\textwidth}}
\toprule
Run & \(\lambda_{\mathrm{video}}\) & \(\lambda_{\mathrm{goal}}\) & \(\lambda_{\mathrm{action}}\) & Outcome & Notes \\
\midrule
D0 & \(0\) & \(0\) & \(0\) & Success 
& Dynamics-only sanity check with only \(\lambda_{\mathrm{dyn}}=10\) active. \\

D1 & \(0\) & \(10\) & \(0\) & Success 
& Goal and soft dynamics active; no action penalty. \\

D2 & \(0\) & \(10\) & \(10^{-3}\) & Success 
& Small action penalty still permits contact-producing actions. \\

D3 & \(0\) & \(10\) & \(10^{-2}\) & Failure 
& Action penalty begins to suppress useful contact actions. \\

D4 & \(0\) & \(10\) & \(5{\times}10^{-2}\) & Failure 
& Actions collapse toward small ineffective motions. \\

D5 & \(1\) & \(10\) & \(10^{-3}\) & Success 
& Adding video loss does not break the successful soft-dynamics configuration. \\
\bottomrule
\end{tabular}
\end{table}

The dynamics-only run succeeds, confirming that the DINO-WM action interface and soft residual are functional. When the task losses are reintroduced, the planner succeeds for \(\lambda_{\mathrm{action}}\in\{0,10^{-3}\}\), but fails for \(10^{-2}\) and \(5\times10^{-2}\), indicating that Push-T requires relatively weak action regularization to preserve contact-producing actions. Re-enabling the video loss with \(\lambda_{\mathrm{action}}=10^{-3}\) also succeeds.

\section{Evaluations Diagnostics}

\label{app:feasibility_eval}
We evaluate the learned feasibility model on the Push-T test set using the checkpoint
\texttt{transformer\_sigma\_delta.pt}. The test set contains \(2440\) samples with macro-actions
\(a_t\in\mathbb{R}^{10}\), corresponding to five packed 2D primitive actions.

\begin{table}[h]
\centering
\small
\caption{Action conversion check from normalized DINO-WM action space to raw environment scale. 
This verifies that the saved action normalization produces plausible Push-T action magnitudes.}
\label{tab:action_conversion_check}
\begin{tabular}{lc}
\toprule
Metric & Value \\
\midrule
Normalized action absolute mean & \(0.8978\) \\
Raw action absolute mean & \(17.9929\) \\
Normalized action global mean & \(-0.1204\) \\
Normalized action global std & \(1.1964\) \\
Raw action global mean & \(-2.5142\) \\
Raw action global std & \(24.0447\) \\
Raw action global min & \(-102.9591\) \\
Raw action global max & \(99.4537\) \\
\bottomrule
\end{tabular}
\end{table}

\begin{table}[h!]
\centering
\small
\caption{DSM latent plausibility evaluation on expert next latents versus small-noise corrupted next latents. 
The DSM energy reliably assigns lower energy to clean expert latents than to corrupted latents.}
\label{tab:dsm_latent_plausibility}
\begin{tabular}{lc}
\toprule
Metric & Value \\
\midrule
Expert energy mean & \(0.2914\) \\
Expert energy median & \(0.2936\) \\
Corrupted energy mean & \(0.3543\) \\
Corrupted energy median & \(0.3567\) \\
Mean margin, corrupted minus expert & \(0.0629\) \\
Fraction expert lower than corrupted & \(1.0000\) \\
Cosine mean, predicted vs. true noise & \(0.5302\) \\
Cosine median, predicted vs. true noise & \(0.5375\) \\
Delta MSE mean & \(0.8569\) \\
Zero-delta MSE mean & \(0.8053\) \\
Delta vs. zero ratio & \(1.0641\) \\
Transition penalty, expert mean & \(0.8569\) \\
Number of samples & \(2440\) \\
\bottomrule
\end{tabular}
\end{table}

\begin{table*}[h]
\centering
\small
\caption{Action discrimination test with fixed expert next latent and different candidate actions. 
The transition head shows weak action discrimination on the test set, while the DSM energy is nearly action-invariant.}
\label{tab:action_discrimination_test}
\begin{tabular}{lcc}
\toprule
Metric & Transition penalty & DSM energy \\
\midrule
Expert action mean & \(0.8569\) & \(0.2914\) \\
Zero action mean & \(0.8530\) & \(0.2912\) \\
Shuffled action mean & \(0.8761\) & \(0.2912\) \\
Negative action mean & \(0.8800\) & \(0.2913\) \\
Gaussian action mean & \(0.8568\) & \(0.2912\) \\
\midrule
Margin: zero minus expert & \(-0.0039\) & -- \\
Margin: shuffled minus expert & \(0.0192\) & -- \\
Margin: negative minus expert & \(0.0231\) & -- \\
Margin: Gaussian minus expert & \(-0.0001\) & -- \\
\midrule
Fraction expert better than zero & \(0.4270\) & \(0.4889\) \\
Fraction expert better than shuffled & \(0.6139\) & \(0.4807\) \\
Fraction expert better than negative & \(0.7115\) & \(0.4869\) \\
Fraction expert better than Gaussian & \(0.5352\) & \(0.4877\) \\
\bottomrule
\end{tabular}
\end{table*}